
\chapter{Background and Related Work}\label{ch:background}

\section{Reinforcement Learning}

\subsection{Markov Decision Processes}

We model the agent-environment interaction as a \emph{Markov Decision Process}
(MDP), defined by the tuple $\mathcal{M} = (\mathcal{S}, \mathcal{A}, P, r, \gamma)$,
where $\mathcal{S}$ is the state space, $\mathcal{A}$ is the action space,
$P: \mathcal{S} \times \mathcal{A} \times \mathcal{S} \to [0,1]$ is the
transition kernel, $r: \mathcal{S} \times \mathcal{A} \to \R$ is the reward
function, and $\gamma \in (0,1)$ is the discount factor.

An agent follows a stochastic policy $\pi_\theta: \mathcal{S} \to
\Delta(\mathcal{A})$, parameterised by $\theta \in \R^d$, where
$\Delta(\mathcal{A})$ denotes the probability simplex over $\mathcal{A}$.
A \emph{trajectory} $\tau = (s_0, a_0, r_0, s_1, a_1, r_1, \ldots, s_T)$ is
generated by executing $\pi_\theta$ in the environment. The agent's objective
is to find the parameters $\theta^*$ that maximise the expected discounted
return:
\[
    J(\theta) = \E*{\sum_{t=0}^{T} \gamma^t r(s_t, a_t)}.
\]

\subsection{Policy Gradient Methods}

The \emph{policy gradient theorem}~\cite{sutton2000policy} provides an
analytical expression for the gradient of $J(\theta)$ with respect to $\theta$:
\[
    \nabla_\theta J(\theta)
    = \E*{\sum_{t=0}^{T} \nabla_\theta \log \pi_\theta(a_t \mid s_t) \cdot G_t},
\]
where $G_t = \sum_{k=t}^{T} \gamma^{k-t} r_k$ is the discounted return from
time step $t$. This identity allows gradient ascent to be performed using
sampled trajectories, without requiring a model of the environment.

The REINFORCE algorithm~\cite{williams1992reinforce} is the canonical Monte
Carlo realisation of this idea. Given a batch of $B$ trajectories sampled from
$\pi_\theta$, the gradient estimate is:
\[
    \hat{g}(\theta) = \frac{1}{B} \sum_{i=1}^{B} \sum_{t=0}^{T}
    \nabla_\theta \log \pi_\theta(a_t^{(i)} \mid s_t^{(i)}) \cdot G_t^{(i)}.
\]
This estimate is unbiased but has high variance, motivating the variance
reduction techniques discussed below.

\subsection{Policy Parameterisation}

For discrete action spaces we use a softmax policy:
\[
    \pi_\theta(a \mid s) = \frac{\exp(f_\theta(s, a))}
    {\sum_{a'} \exp(f_\theta(s, a'))},
\]
where $f_\theta$ is a multilayer perceptron (MLP) with parameters $\theta$.
For continuous action spaces we use a Gaussian policy whose mean is output by
the network and whose variance is a learned parameter.

\section{Federated Learning}

\subsection{The Federated Averaging Setting}

Federated learning~\cite{mcmahan2017fedavg} considers $K$ workers, each holding
a local dataset, that collaboratively train a shared model without sharing their
data. In each round, the server broadcasts the current model to all workers.
Each worker computes a local gradient update and sends it back to the server.
The server aggregates the updates --- typically by a weighted average --- and
applies the result to the global model.

In the FedRL setting, the ``local dataset'' is replaced by a set of trajectories
sampled by the worker from its local environment instance. The locally computed
quantity is a policy gradient estimate rather than a supervised loss gradient.
The server aggregates these gradient estimates and updates the shared policy
parameters.

\subsection{Challenges Specific to Federated RL}

Federated RL inherits the standard challenges of federated learning (communication
overhead, statistical heterogeneity across workers) and adds several specific
difficulties:

\begin{itemize}
    \item \textbf{Non-stationarity.} The distribution of trajectories depends on
          the current policy. As the policy improves, the gradient statistics change,
          requiring aggregation methods that are robust to non-stationary inputs.
    \item \textbf{Variable trajectory length.} Unlike supervised learning minibatches,
          trajectories have variable length determined by the episode dynamics. This
          makes communication overhead difficult to bound a priori.
    \item \textbf{High gradient variance.} Policy gradient estimates have
          notoriously high variance even in the centralised setting. In the federated
          setting, variance from environment stochasticity compounds with variance
          from the heterogeneous environments of different workers.
\end{itemize}

\section{Byzantine Robustness}

\subsection{The Byzantine Generals Problem}

The \emph{Byzantine fault model}~\cite{lamport1982byzantine} considers a system
where an unknown subset of $f$ out of $K$ participants may behave arbitrarily ---
sending corrupted, manipulated, or strategically chosen messages to the
aggregation server. In the federated learning context, Byzantine workers send
malicious gradient updates designed to corrupt the shared model.

The challenge is that the server cannot distinguish a malicious update from a
legitimate one based on its value alone, since honest workers may produce
high-variance gradients due to environment stochasticity. Robust aggregation
rules must therefore be designed to tolerate a bounded fraction of Byzantine
workers while still incorporating useful information from the honest majority.

\subsection{Gradient-Based Byzantine Filters}

Several aggregation rules provide Byzantine robustness by inspecting the
statistical properties of the submitted gradients. Coordinate-wise median and
trimmed mean are classical examples. More specialised filters exploit the
geometry of gradient vectors: if $f$ workers are Byzantine and $K - f$ workers
are honest, a filter based on pairwise distances can exclude outliers.

FedPG-BR uses a \emph{probabilistic} filter based on the empirical distribution
of gradient norms, parameterised by $\sigma$ and $\delta$. Crucially, this
filter is calibrated to the expected gradient statistics under the current
policy, which means its behaviour changes as the policy evolves --- an
interaction we examine in detail in \Cref{ch:system}.

\section{FedPG-BR}\label{sec:fedpgbr}

FedPG-BR~\cite{flint2023fedpgbr} extends the federated policy gradient
framework with two key algorithmic contributions: a Byzantine filter and a
variance reduction technique based on SCSG.

\subsection{Byzantine Filtering}

Let $g_1, \ldots, g_K \in \R^d$ denote the gradient vectors submitted by the
$K$ workers. FedPG-BR computes the empirical mean $\bar{g}$ and uses a
probabilistic norm-based criterion to identify and exclude workers whose
gradient deviates excessively from the group:
\[
    \mathcal{G} = \left\{ i : \norm*{g_i - \bar{g}} \le \sigma \cdot \sqrt{d} \right\}.
\]
The threshold $\sigma$ controls the sensitivity of the filter. A small $\sigma$
filters aggressively; a large $\sigma$ is permissive. The aggregated gradient
is the mean over the good agents $\mathcal{G}$:
\[
    \mu_t = \frac{1}{|\mathcal{G}|} \sum_{i \in \mathcal{G}} g_i.
\]

\subsection{SCSG Variance Reduction}

To reduce the variance of the policy gradient estimate, FedPG-BR uses the
\emph{Stochastic Controlled Stochastic Gradient} (SCSG) technique, also
explored independently in SVRPG~\cite{papini2018svrpg}. Starting from the
aggregated gradient $\mu_t$ at policy $\theta_t^0$, the server performs a
random number $N_t \sim \mathrm{Geom}(p)$ of variance-reduced inner steps.
Each inner step uses importance sampling to compute a variance-reduced
gradient $v_{t,n}$:
\[
    v_{t,n} = \hat{g}(\theta_t^n) - \hat{g}_{\text{old}}(\theta_t^n) + \mu_t,
\]
where $\hat{g}(\theta_t^n)$ is a mini-batch gradient at the current parameters
and $\hat{g}_{\text{old}}(\theta_t^n)$ is a corrected gradient from the
snapshot policy $\theta_t^0$, weighted by the importance ratio:
\[
    \rho = \frac{\pi_{\theta_t^0}(a \mid s)}{\pi_{\theta_t^n}(a \mid s)}.
\]
If the importance ratio drifts outside the window $[0.995, 1.005]$, the inner
loop terminates to avoid excessive bias from off-policy corrections.

\section{Related Work}

\subsection{Federated Policy Gradient Algorithms}

The federated policy gradient literature has grown rapidly. \cite{papini2018svrpg}
introduced SVRPG, which applies SCSG variance reduction to the centralised
policy gradient setting and forms the basis for the variance reduction component
in FedPG-BR. Several works have studied convergence guarantees for federated
policy gradient under various assumptions on the environment heterogeneity and
Byzantine fraction.

\subsection{Benchmarking in Federated Learning}

The supervised federated learning community has developed several benchmarking
frameworks, including Flower~\cite{beutel2020flower}, LEAF, and FedML. These
frameworks standardise the communication layer and allow algorithms to be
compared under controlled conditions. The Flower framework in particular has
gained wide adoption for its clean Python API and support for both simulation
and real distributed deployment.

No analogous benchmark exists for federated reinforcement learning. This gap
motivates the present work.

\subsection{Gymnasium and MuJoCo}

Environments are drawn from the Gymnasium library~\cite{brockman2016gym}, the
community-maintained successor to OpenAI Gym. Continuous control experiments
use MuJoCo~\cite{todorov2012mujoco}, a physics engine widely used in RL
research. The specific environments used are CartPole-v1, Acrobot-v1,
MountainCar-v0, LunarLander-v3, and HalfCheetah-v5.
